% Schematic for the backward chain in Lemma 6.3.
% The highlighted vertical intervals are XO patches; the double-headed links
% are precisely the successful interactions used by the chain.
\begingroup
\definecolor{convtrailred}{RGB}{175,32,38}
\definecolor{convtrailoutline}{gray}{0.18}
\definecolor{convtrailcone}{gray}{0.88}
\definecolor{convtrailconeborder}{gray}{0.25}
\definecolor{convtrailsecondary}{gray}{0.43}
\begin{tikzpicture}[x=0.56cm,y=0.56cm]
\path[use as bounding box] (0.20,-0.12) rectangle (7.22,7.48);
\tikzset{
  convtrail cone/.style={fill=convtrailcone,draw=convtrailconeborder,
    line width=0.62pt,line join=round},
  convtrail patch/.style={draw=convtrailred,line width=2.65pt,
    line cap=round},
  convtrail interaction/.style={draw=convtrailoutline,line width=0.78pt,
    shorten <=1.2pt,shorten >=1.2pt,
    {Stealth[length=1.35mm,width=0.96mm]}-{Stealth[length=1.35mm,width=0.96mm]}},
  convtrail axis/.style={draw=convtrailoutline,line width=0.62pt,
    -{Stealth[length=1.8mm,width=1.2mm]}},
  convtrail tick/.style={draw=convtrailoutline,line width=0.52pt},
  convtrail baseline/.style={draw=convtrailoutline,line width=0.52pt},
  convtrail guide/.style={draw=convtrailsecondary,line width=0.48pt,
    dash pattern=on 2.2pt off 2.0pt},
  convtrail label/.style={font=\scriptsize,text=convtrailoutline,inner sep=1pt},
  convtrail secondary label/.style={font=\scriptsize,
    text=convtrailsecondary,inner sep=1pt}
}

% A deliberately irregular spacetime outline of the interaction cone.
\path[convtrail cone]
  (3.26,0)
  .. controls (3.07,0.55) and (2.95,1.08) .. (2.98,1.56)
  .. controls (2.99,1.91) and (2.45,2.17) .. (2.28,2.63)
  .. controls (2.10,3.11) and (2.48,3.48) .. (2.25,3.87)
  .. controls (2.00,4.31) and (1.45,4.50) .. (1.30,5.03)
  .. controls (1.15,5.55) and (0.82,6.01) .. (0.66,6.50)
  -- (5.56,6.50)
  .. controls (5.44,6.04) and (5.06,5.68) .. (5.25,5.18)
  .. controls (5.43,4.70) and (5.06,4.43) .. (4.87,4.08)
  .. controls (4.64,3.65) and (5.07,3.31) .. (4.92,2.87)
  .. controls (4.77,2.40) and (4.34,2.04) .. (4.47,1.58)
  .. controls (4.60,1.11) and (4.12,0.54) .. (3.95,0)
  -- cycle;

\node[convtrail secondary label,anchor=west] at (1.38,5.62)
  {$\mb{Cone}$};

% The backward chain.  Its patch lifetimes partition [0,u).
\draw[convtrail baseline] (0.40,0)--(5.78,0);
\draw[convtrail patch] (3.35,0)--(3.35,1.34);
\draw[convtrail interaction] (3.35,1.34)--(4.45,1.34);
\draw[convtrail patch] (4.45,1.34)--(4.45,2.66);
\draw[convtrail interaction] (2.75,2.66)--(4.45,2.66);
\draw[convtrail patch] (2.75,2.66)--(2.75,3.98);
\draw[convtrail interaction] (2.75,3.98)--(4.55,3.98);
\draw[convtrail patch] (4.55,3.98)--(4.55,5.18);
\draw[convtrail interaction] (3.25,5.18)--(4.55,5.18);
\draw[convtrail patch] (3.25,5.18)--(3.25,6.50);

% The prescribed successful interaction at time u.
\draw[convtrail interaction] (3.25,6.50)--(5.12,6.50);
\node[convtrail label,anchor=south] at (4.18,6.66) {$(i,u,S)$};

% Time axis and the guide locating the final interaction time.
\draw[convtrail guide] (5.12,6.50)--(6.38,6.50);
\draw[convtrail axis] (6.38,0)--(6.38,7.28);
\draw[convtrail tick] (6.28,0)--(6.48,0);
\draw[convtrail tick] (6.28,6.50)--(6.48,6.50);
\node[convtrail label,anchor=west] at (6.54,0) {$0$};
\node[convtrail label,anchor=west] at (6.54,6.50) {$u$};
\node[convtrail secondary label,anchor=south] at (6.38,7.29)
  {time};
\end{tikzpicture}
\endgroup
